% ========================================================
% This document is a customizable CV/Resume template built using LaTeX.
% The template is designed for easy customization and clear structure.
%
% Author: Matthew DeVerna (www.matthewdeverna.com)
% Date: 2024
% Design: Cased on Hause Lin's CV (hauselin.com)
%
% Project Overview:
% -----------------
% This LaTeX document is designed to help you create a professional CV or
% resume with ease. It uses as little fancy LaTex functionality or custom functions as possible to maximize its longterm durability and flexibility.
% The document is structured into multiple sections, each loaded from
% separate subfiles for modularity and ease of maintenance.
%
% Key Features:
% -------------
% - Customizable sections: Education, Research Experience, Awards, Publications, etc.
% - Bookmarks in the PDF for easy navigation
% - Styled bibliography with BibLaTeX
% - Hyperlinked email and website
% - FontAwesome icons for additional styling
%
% Getting Started:
% ----------------
% 1. Customize your personal information by modifying the \mytitle command.
% 2. Add your content to the respective subfiles (e.g., education.tex, exp_research.tex).
% 3. Update the bibliography file (ref.bib) with your publications and categorize them with keywords.
%
% Important Notes:
% ----------------
% - This main file includes the overall structure and settings.
% - Each section has its own detailed instructions for further customization.
%
% - The sections included below mirror what the website currently shows. The
%   template also ships with tools.tex, exp_other.tex, advising.tex and
%   service.tex, and with slots for journal articles, working papers, talks,
%   posters and demonstrations. Those are left out rather than filled with
%   placeholder content; add the section back here when there is something
%   real to put in it.
%
% ========================================================

\documentclass[11pt]{article} % Choose the document class and font size
\usepackage[margin=1in]{geometry} % Page layout settings

% Set the citation style
\usepackage[
    backend=biber,      % Specifies the backend to be used by BibLaTeX for processing the bibliography. 'biber' is the default backend.
    maxnames=20,        % Limits the maximum number of author names to display before abbreviating with "et al."
    style=nature,       % Sets the citation style to 'nature,' which is commonly used in scientific papers.
    sorting=ydnt,       % Specifies the sorting order of entries in the bibliography:
                        % y - year (descending)
                        % d - descending order
                        % n - name
                        % t - title
    defernumbers=true,  % Delays the assignment of citation numbers until the end of the document, allowing for the correct order of citations within each bibliography section.
]{biblatex}
\addbibresource{ref.bib} % Adds the bibliography resource file 'ref.bib' containing all the references.

% Allows columns that stretch across pages
\usepackage{longtable}

% Table functionality and beautification (not strictly needed)
\usepackage{bookmark}

% Use icons, if you want.
% All available icons: http://mirrors.ibiblio.org/CTAN/fonts/fontawesome5/doc/fontawesome5.pdf
\usepackage{fontawesome}

% Allows font justification control (needed for clean pub-list formatting)
\usepackage{ragged2e}

% For underlining with line breaks
\usepackage{soul}

% All fonts: https://tug.org/FontCatalogue/
\usepackage{kpfonts} % More professional font
% \usepackage[default]{sourcecodepro} % Code-like font
\usepackage[T1]{fontenc}

% Control hyperlinks and colors
% CUSTOM COLORS INCLUDED DIRECTLY AFTER \begin{document}
\usepackage{xcolor}
\usepackage{hyperref}
\hypersetup{
    colorlinks=true,        % Enable colored links
    breaklinks=true,        % Allow links to break across lines
    linkcolor=cornflowerblue,    % Color of internal links
    urlcolor=cornflowerblue,     % Color of URL links
    anchorcolor=cornflowerblue,  % Color of anchors
    citecolor=cornflowerblue,    % Color of citations
    pdftitle={Keito Kozaki -- Curriculum Vitae},    % Title of the PDF
    pdfauthor={Keito Kozaki}, % Author of the PDF
    bookmarksopen=true,      % Open bookmarks panel at start
}

%%% CONVENIENCE FUNCTIONS GO HERE %%%
%%% ----------------------------- %%%
\newcommand{\mytitle}[4]{
  \begin{center}
    \Large\textbf{#1}\normalsize \\ % Name in large bold font
    \href{mailto:#2}{#2} \\ % Email with mailto: link
    \href{https://#3}{#3} \\ % Website with link
    #4 % Address
  \end{center}
}
%%% ----------------------------- %%%


\begin{document}
% Set custom colors here (imported directly after \begin{document})
% The below use HTML hex codes.
% More HTML hex codes: https://encycolorpedia.com/html
\definecolor{firebrick}{HTML}{b22222} 
\definecolor{darkslategrey}{HTML}{2f4f4f} 
\definecolor{cornflowerblue}{HTML}{6495ed} 
\definecolor{mediumslateblue}{HTML}{7b68ee}  % Load custom colors from colors file
\mytitle
    {Keito Kozaki}
    {kozaki@lmd.ist.hokudai.ac.jp}
    {keito0329.github.io}
    {Media Dynamics Laboratory, Graduate School of Information Science and Technology\\
     Hokkaido University, Sapporo, Japan\\
     \href{https://orcid.org/0009-0004-5149-1249}{ORCID: 0009-0004-5149-1249}}


% Ensure right side margin is not surpassed by bibliography and the right margin is aligned throughout
\RaggedRight


% These \pdfbookmark lines create bookmarks in the exported PDF document that display in the left pane.
% Value in [] sets the indentation level of the bookmark
\pdfbookmark[1]{Education}{}
\section*{Education}
% Add your educational background here!

\begin{longtable}[l]{@{}p{.125\textwidth} p{0.875\textwidth}}

    2026-- & \textbf{M.S.}, Information Science and Technology, Hokkaido University, \href{https://www-lmd.ist.hokudai.ac.jp/}{Media Dynamics Laboratory} \\

    2022--26 & \textbf{B.Eng.}, School of Engineering, Hokkaido University

\end{longtable}


\pdfbookmark[1]{Research Experience}{exp_research}
\section*{Research Experience}
\label{exp_research}
% List your research experience here

\begin{longtable}[l]{@{}p{.125\textwidth} p{0.875\textwidth}}

    2026-- & Research Assistant, \href{https://www-lmd.ist.hokudai.ac.jp/}{Media Dynamics Laboratory}, Hokkaido University \\

\end{longtable}


\pdfbookmark[1]{Industry Experience}{exp_industry}
\section*{Industry Experience}
\label{exp_industry}
% Positions held outside the university. The template also ships exp_other.tex
% for experience that does not fit anywhere else; this is kept separate so the
% section can be named for what it actually is.

\begin{longtable}[l]{@{}p{.125\textwidth} p{0.875\textwidth}}

    2026-- & Applied Engineer (part-time), Simulacra Inc. \\

\end{longtable}


\pdfbookmark[1]{Awards \& Honors}{awards}
\section*{Awards \& Honors}
\label{awards}
% Include awards, honors, grant funding, etc. here

\begin{longtable}[l]{@{}p{.125\textwidth} p{0.875\textwidth}}

    2026 & Dean's Award, School of Engineering, Hokkaido University \\

    2026 & Student Presentation Award, DEIM2026 \\

    2025 & Young Researcher Excellent Paper Presentation Award, Joint Conference of Hokkaido Chapters of Electrical and Information Institutes \\

    2025 & Excellent Paper Award, Silver Prize, IEEE GCCE 2025 \\

\end{longtable}


\pdfbookmark[1]{Scholarships \& Fellowships}{scholarships}
\section*{Scholarships \& Fellowships}
\label{scholarships}
% Scholarships and fellowships. Kept apart from awards.tex because these are
% funding rather than prizes, matching how the website lists them.

\begin{longtable}[l]{@{}p{.125\textwidth} p{0.875\textwidth}}

    2026-- & \href{https://tcs-foundation.or.jp/}{TCS Scholarship} \\

\end{longtable}



\pdfbookmark[1]{Publications}{pubs}
\section*{Publications}
\label{pubs}

% Add a Google Scholar link here once the profile exists, following the
% template's original form:
% \vspace{-.75em}
% \small
% \faGoogle~\href{https://scholar.google.com/citations?user=XXXX}{Google Scholar}
% \normalsize

% Ensures publications which are not cited in the document are included in the above sections
\nocite{*} % Ensures uncited items are included

\pdfbookmark[2]{International Conference Papers}{conferences}
\subsection*{International Conference Papers}
\label{conferences}
\newrefcontext[labelprefix=C] % Will prefix bibliography numbers with this letter
\printbibliography[
    type=inproceedings, % Only include @inproceedings ref.bib items
    heading=none,       % Do not include header. Gives us more control.
    resetnumbers=true,  % Start item counter from zero
    keyword=IC          % Include items in ref.bib with keyword={IC}
]

\pdfbookmark[2]{Domestic Conference Papers}{domestic}
\subsection*{Domestic Conference Papers}
\label{domestic}
\newrefcontext[labelprefix=D]
\printbibliography[type=inproceedings,heading=none,resetnumbers=true,keyword=DC]


\pdfbookmark[1]{Selected Media Coverage}{media}
\section*{Selected Media Coverage}
\label{media}
%  Include media coverage here
%
% Entries are written in English because this document is built with pdfLaTeX,
% which cannot typeset the Japanese titles; items published in Japanese say so.
%
\begin{longtable}[l]{@{}p{.125\textwidth} p{0.875\textwidth}}

    2026 & Hokkaido University, \href{https://www.hokudai.ac.jp/news/2026/10/ai-14.html}{press release on the RecSys 2026 paper} -- In Japanese. \\

    2026 & Top Information Retrieval Papers of the Week, \href{https://recsys.substack.com/p/do-latent-reasoning-steps-actually}{Do Latent Reasoning Steps Actually Help Retrieval?, Adapting Dense Retrievers from Agent Search Traces, and More!} \\

\end{longtable}
 % Input lines load the material from the subdocuments


\pdfbookmark[1]{Teaching}{teaching}
\section*{Teaching}
\label{teaching}

\subsection*{Hokkaido University}
%
\begin{longtable}[l]{@{}p{.125\textwidth} p{0.875\textwidth}}

    2026 & Teaching Assistant, Exercise in Media Network II \\

    2026 & Teaching Assistant, Media Network Laboratory IA \\

\end{longtable}



% Pretty ending with the date last updated
\centering
\rule{0.25\linewidth}{0.4pt}\\
\medskip
Last updated: \today

\end{document}
